\chapter{Containers, Remote Execution, and HPC Schedulers}
\label{ch:backends}

Part~II extends the foundations of Part~I across the engine's heterogeneous
execution backends. This chapter covers Lessons~8 and~9: Docker-backed
execution, remote SSH execution, and HPC job-scheduler (Slurm) execution ---
the three backends most directly associated with classical HPC practice.

\begin{lessonbox}{Lesson 8 --- Run Docker-Backed Tasks}
\textbf{Objective.} Run a task inside a Docker container.

\textbf{Concept.} \texttt{TaskType.DOCKER} creates a \texttt{DockerTask}. The
task uses the Docker Python SDK, pulls or uses an image, starts a container,
and runs the task command with \texttt{docker exec}.

\textbf{Prerequisites.} Docker installed; Docker daemon running; current user
allowed to access Docker. Verify with \texttt{docker info}.

\textbf{Code} (\texttt{lesson\_08.py}):
\begin{lstlisting}
from dagon import Workflow
from dagon.task import DagonTask, TaskType


if __name__ == "__main__":
    workflow = Workflow("Lesson08")

    task = DagonTask(
        TaskType.DOCKER,
        "container_python",
        "python --version > python_version.txt",
        image="python:3.12-slim",
    )

    workflow.add_task(task)
    workflow.make_dependencies()
    workflow.run()
\end{lstlisting}

\textbf{Verify:}
\begin{lstlisting}[style=bash]
DIR=$(find /tmp -maxdepth 1 -type d -name '*-container_python' | tail -n 1)
sed -n '1,80p' "$DIR/python_version.txt"
\end{lstlisting}

\textbf{Expected behavior.} The output file contains the Python version
reported inside the container.

\textbf{Scientific practice note.} Containers help preserve software
environments. For published workflows, record the exact image tag or digest
used for scientific results.
\end{lessonbox}

\section{Discussion: Containers as Environment Provenance}

Lesson~8 illustrates that, architecturally, a \texttt{DockerTask} differs
from a \texttt{Batch} task only in \emph{where} the launcher executes ---
inside a container via \texttt{docker exec} rather than directly via
\texttt{subprocess.Popen}. Every dependency, staging, and checkpoint
mechanism developed in Part~I applies unchanged. The scientific value is
environment provenance: \texttt{image="python:3.12-slim"} is, in a published
workflow, as much a part of the reproducible record as the command string
itself, and should be pinned to a digest rather than a mutable tag whenever a
result is meant to be exactly reproducible.

\begin{lessonbox}{Lesson 9 --- Prepare Remote and Slurm Workflows}
\textbf{Objective.} Understand how to construct remote and Slurm tasks, and
how to verify the environment before running them.

\textbf{Concept.} Remote task variants are selected by constructor arguments
such as \texttt{ip}. Slurm tasks use \texttt{TaskType.SLURM} and generate
\texttt{sbatch} commands.

\textbf{Remote batch task:}
\begin{lstlisting}
from dagon import Workflow
from dagon.task import DagonTask, TaskType


workflow = Workflow("RemoteExample")

remote = DagonTask(
    TaskType.BATCH,
    "remote_hostname",
    "hostname > host.txt",
    ip="example.org",
    ssh_username="researcher",
    keypath="/home/researcher/.ssh/id_rsa",
)

workflow.add_task(remote)
workflow.make_dependencies()
workflow.run()
\end{lstlisting}

Before running, verify SSH manually:
\begin{lstlisting}[style=bash]
ssh -i /home/researcher/.ssh/id_rsa researcher@example.org hostname
\end{lstlisting}

\textbf{Slurm task:}
\begin{lstlisting}
from dagon import Workflow
from dagon.task import DagonTask, TaskType


workflow = Workflow("SlurmExample")

task = DagonTask(
    TaskType.SLURM,
    "slurm_hostname",
    "hostname > host.txt",
    partition="compute",
    ntasks=1,
    time="00:05:00",
)

workflow.add_task(task)
workflow.make_dependencies()
workflow.run()
\end{lstlisting}

Before running, verify Slurm:
\begin{lstlisting}[style=bash]
which sbatch
sinfo
\end{lstlisting}

\textbf{Local structural verification} (no scheduler required):
\begin{lstlisting}
from dagon.task import DagonTask, TaskType

task = DagonTask(TaskType.SLURM, "check", "echo check", partition="debug", ntasks=1, working_dir="scratch")
print(type(task).__name__)
print(task.generate_command("launcher.sh"))
\end{lstlisting}

\textbf{Scientific practice note.} Remote and HPC workflows are reproducible
only when resource requests, environment modules, host configuration, and
input data locations are recorded.
\end{lessonbox}

\section{Discussion: The Factory Pattern for Remote Variants}

Lesson~9 makes explicit a design decision already anticipated in
\Cref{ch:architecture}: \emph{remote} execution is not a separate
\texttt{TaskType}, but a keyword-argument-triggered specialization of an
existing one. Passing \texttt{ip}, \texttt{ssh\_username}, and optionally
\texttt{keypath} to a \texttt{TaskType.BATCH} construction causes
\texttt{Batch.\_\_new\_\_} to return a \texttt{RemoteBatch} instance instead
(\Cref{ch:architecture}); the same pattern extends to remote Slurm and remote
Docker. This keeps the DAG-authoring surface uniform --- a workflow author
reasons about \emph{tasks and their dependencies}, not about a proliferating
taxonomy of remote/local task classes.

\begin{table}[h]
\centering
\caption{Backend selection by scientific software need (after the project's
own guidance).}
\label{tab:backend-selection}
\begin{tabular}{@{}ll@{}}
\toprule
\textbf{Need} & \textbf{Typical backend} \\
\midrule
Small local utility & \texttt{BATCH} \\
Compiled HPC model & \texttt{SLURM} \\
Dependency-heavy tool & \texttt{DOCKER} \\
Remote instrument/service host & remote \texttt{BATCH} \\
Elastic compute & \texttt{CLOUD} \\
\bottomrule
\end{tabular}
\end{table}

\section{The Local-Structural-Verification Pattern}

A recurring theme across the remote, Slurm, and (later) meta-workflow lessons
is that a workshop laptop or CI environment frequently lacks the external
service (a reachable SSH host, a live Slurm controller) needed for full
execution. Lesson~9's structural-verification snippet --- constructing the
task object and calling \texttt{generate\_command()} without ever calling
\texttt{workflow.run()} --- is the project's recommended substitute: it
exercises the \emph{command-generation} logic exhaustively while requiring no
external service at all, and is the pattern this book recommends whenever a
reader wants to validate a Slurm or remote task specification before
deploying it to a real cluster.

\section{Containers and the Reproducibility Problem, Revisited}

\Cref{ch:backends}'s discussion of Lesson~8 already noted that container
image identifiers are part of a workflow's scientific record. This deserves
expansion, because the distinction between a mutable \emph{tag} and an
immutable \emph{digest} is one of the more consequential, and most
frequently overlooked, reproducibility decisions a workflow author using
\texttt{DOCKER} tasks will make.

A tag such as \texttt{python:3.12-slim} is a mutable pointer: the
maintainers of that image may, and routinely do, push a new image under the
same tag --- to apply a security patch, update a base operating system, or
bump a minor library version --- without any change to the tag string
itself. A workflow that pins only the tag is therefore not fully
reproducible in the strict sense this book has used throughout: rerunning
the identical workflow source code six months later may silently execute a
different, updated image, with no error, warning, or visible change to the
workflow definition. A \emph{digest}, by contrast --- a cryptographic hash
of the image's content, typically written as
\texttt{image@sha256:<hex-digest>} --- refers to one, and only ever one,
immutable image content. Pinning to a digest sacrifices automatic access to
upstream patches but gains the property that the phrase ``the same
container image'' becomes mathematically, not merely conventionally, true
across arbitrarily long timespans.

\begin{practicebox}
For workflows whose results are destined for publication, archival, or
long-term operational reuse --- the environmental forecasting deployments
discussed in \Cref{ch:case-study} being the clearest example --- this book
recommends recording, alongside the workflow source, both the human-readable
tag (for legibility) and the resolved digest (for strict reproducibility) at
the time a result was produced, even though the \texttt{DockerTask}
constructor itself accepts only the \texttt{image} string and does not
distinguish the two forms syntactically.
\end{practicebox}

\section{The Economics of Remote and Scheduled Execution}

Lesson~9's two backends --- remote SSH execution and Slurm scheduling ---
share a structural property worth making explicit: both introduce a
\emph{queueing} or \emph{latency} cost that purely local \texttt{BATCH}
execution does not incur, and a workflow author reasoning about overall
pipeline turnaround time should account for this cost as a first-class
design consideration, not an incidental detail.

A local \texttt{BATCH} task's total wall-clock cost is, to a first
approximation, just the time the underlying command itself takes to run,
plus the comparatively small overhead of scratch-directory creation and
staging (\Cref{ch:lessons-03-04}). A \texttt{SLURM} task's total wall-clock
cost, by contrast, is the sum of the command's own runtime \emph{and} the
time the submitted job spends queued, waiting for the scheduler to allocate
the requested resources --- a quantity that depends on cluster-wide
contention entirely outside the workflow's own control, and that can, on a
busy shared HPC system, dwarf the actual computational time by a large
factor. A remote \texttt{BATCH} task incurs a smaller, but still non-zero,
analogous cost: SSH connection establishment and, depending on network
conditions, command round-trip latency.

\begin{practicebox}
This has a direct consequence for how the asynchronous launch mechanism
developed in \Cref{ch:meta-async} should be read in anticipation: a
workflow mixing local and Slurm-scheduled tasks benefits disproportionately
from \texttt{launch()}/\texttt{wait()} non-blocking execution precisely
\emph{because} Slurm queueing time is unpredictable and potentially long ---
a controller blocked synchronously on \texttt{run()} during a multi-hour
queue wait cannot do anything else useful in the meantime, whereas a
controller using \texttt{launch()} remains free to handle other requests,
report progress, or even submit additional, independent workflows, while
the Slurm-scheduled portion of the DAG waits its turn.
\end{practicebox}

\section{A Worked Comparison Across the Backend Taxonomy}

To consolidate \Cref{tab:backend-selection}'s qualitative guidance into a
more concrete worked illustration, consider a single post-processing
operation --- regridding a model's native output onto a regular
latitude-longitude grid --- and how the choice of backend would change the
shape of the corresponding \texttt{DagonTask} construction, without changing
the regridding logic itself.

As a local utility, assuming the regridding tool is already installed on the
machine running the workflow controller, the task is simply:
\begin{lstlisting}
DagonTask(TaskType.BATCH, "regrid", "regrid_tool --in workflow:///model/out.nc --out regridded.nc")
\end{lstlisting}

As a dependency-heavy tool requiring a specific, hard-to-install geospatial
library stack, the same logical operation becomes a \texttt{DOCKER} task,
trading installation complexity for container-image maintenance:
\begin{lstlisting}
DagonTask(TaskType.DOCKER, "regrid", "regrid_tool --in workflow:///model/out.nc --out regridded.nc", image="lab/regrid-tools:2.3")
\end{lstlisting}

As a compute-intensive operation on a very large grid, justifying scheduler-
managed resource allocation, it becomes a \texttt{SLURM} task, trading local
immediacy for queued, resource-guaranteed execution:
\begin{lstlisting}
DagonTask(TaskType.SLURM, "regrid", "regrid_tool --in workflow:///model/out.nc --out regridded.nc", partition="compute", ntasks=8, time="00:30:00")
\end{lstlisting}

\begin{practicebox}
Note what does \emph{not} change across all three variants: the
\texttt{workflow:///model/out.nc} dependency reference, and therefore the
dependency edge it induces, the staging behavior that places the file in the
task's working directory, and the task's name within the DAG. This
invariance is the architectural payoff promised in \Cref{ch:architecture}'s
discussion of the factory pattern (\Cref{sec:factory-pattern-depth}):
changing \emph{where} a logical operation executes is, in the common case, a
matter of changing the \texttt{TaskType} and a small number of backend-
specific keyword arguments, not of redesigning the workflow's dependency
structure.
\end{practicebox}

\section{Chapter Summary}

Docker, remote-batch, and Slurm execution share one architectural lesson:
heterogeneity of \emph{where} a command runs is isolated entirely within the
concrete task subclass's execution hook, while dependency resolution,
staging, checkpointing, and the \texttt{workflow://} schema are identical to
the purely local workflows of Part~I. This chapter additionally examined the
tag-versus-digest reproducibility distinction for container images, reasoned
about the queueing and latency costs that distinguish remote and
scheduler-backed execution from purely local execution, and worked through a
single logical operation expressed across three different backends to make
concrete the claim that backend choice is, architecturally, orthogonal to
dependency-graph structure. \Cref{ch:meta-async} now turns to composing
multiple workflows together and to controlling \emph{when}, rather than
merely \emph{where}, execution happens.

\begin{exercise}
Using the worked regridding example above as a template, write the
\texttt{TaskType.BATCH} \emph{remote} variant of the same task (executing on
a named remote host via SSH, as in Lesson~9), and explain, referencing
\Cref{sec:factory-pattern-depth}, exactly which class \texttt{DagonTask}
would instantiate for this fourth variant, and how that class relates to
the plain local \texttt{Batch} class used in the first variant above.
\end{exercise}
